\documentclass[10pt,a4paper]{amsart}
\usepackage{amsmath,amssymb,amsthm,graphicx}
% Retained preamble for the Stacks Project chapter snapshot in this directory.
% The chapter's own preamble.tex is a copy of the shared Stacks Project preamble; the block
% below is sliced MECHANICALLY, line for line, from it. Nothing was added, removed or
% redefined; the wrapper class, the xr-hyper/hyperref wiring and the externaldocument list
% are omitted because the assembler supplies the document class.
%
% lines taken from preamble.tex: 3, 5, 9, 11, 12, 156, 157, 158, ... (51 lines total)
%
% package coverage restored from the chapter's own preamble (2026-10-13): xy, xy-2cell

\usepackage[all]{xy}
\xyoption{2cell}
\UseAllTwocells

\newenvironment{reference}{\comment}{\endcomment}
\newenvironment{slogan}{\comment}{\endcomment}
\newenvironment{history}{\comment}{\endcomment}
\theoremstyle{plain}
\newtheorem{theorem}[subsection]{Theorem}
\newtheorem{proposition}[subsection]{Proposition}
\newtheorem{lemma}[subsection]{Lemma}
\theoremstyle{definition}
\newtheorem{definition}[subsection]{Definition}
\newtheorem{example}[subsection]{Example}
\newtheorem{exercise}[subsection]{Exercise}
\newtheorem{situation}[subsection]{Situation}
\theoremstyle{remark}
\newtheorem{remark}[subsection]{Remark}
\newtheorem{remarks}[subsection]{Remarks}
\numberwithin{equation}{subsection}
\def\lim{\mathop{\mathrm{lim}}\nolimits}
\def\colim{\mathop{\mathrm{colim}}\nolimits}
\def\Spec{\mathop{\mathrm{Spec}}}
\def\Hom{\mathop{\mathrm{Hom}}\nolimits}
\def\Ext{\mathop{\mathrm{Ext}}\nolimits}
\def\SheafHom{\mathop{\mathcal{H}\!\mathit{om}}\nolimits}
\def\SheafExt{\mathop{\mathcal{E}\!\mathit{xt}}\nolimits}
\def\Sch{\mathit{Sch}}
\def\Mor{\mathop{\mathrm{Mor}}\nolimits}
\def\Ob{\mathop{\mathrm{Ob}}\nolimits}
\def\Sh{\mathop{\mathit{Sh}}\nolimits}
\def\NL{\mathop{N\!L}\nolimits}
\def\CH{\mathop{\mathrm{CH}}\nolimits}
\def\proetale{{pro\text{-}\acute{e}tale}}
\def\etale{{\acute{e}tale}}
\def\QCoh{\mathit{QCoh}}
\def\Ker{\mathop{\mathrm{Ker}}}
\def\Im{\mathop{\mathrm{Im}}}
\def\Coker{\mathop{\mathrm{Coker}}}
\def\Coim{\mathop{\mathrm{Coim}}}
\DeclareMathSymbol{\boxtimes}{\mathbin}{AMSa}{"02}
\def\QCohstack{\mathcal{QC}\!\mathit{oh}}
\def\Cohstack{\mathcal{C}\!\mathit{oh}}
\def\Spacesstack{\mathcal{S}\!\mathit{paces}}
\def\Quotfunctor{\mathrm{Quot}}
\def\Hilbfunctor{\mathrm{Hilb}}
\def\Curvesstack{\mathcal{C}\!\mathit{urves}}
\def\Polarizedstack{\mathcal{P}\!\mathit{olarized}}
\def\Complexesstack{\mathcal{C}\!\mathit{omplexes}}
\def\Pic{\mathop{\mathrm{Pic}}\nolimits}
\def\Picardstack{\mathcal{P}\!\mathit{ic}}
\def\Picardfunctor{\mathrm{Pic}}
\def\Deformationcategory{\mathcal{D}\!\mathit{ef}}
\title{Flatness modulo an ideal plus vanishing of Tor one implies flatness on nilpotently annihilated modules}
\author{The Stacks Project Authors}
\date{stacks-project:algebra.tex:local-snapshot-20261012; GNU FDL 1.2 or later}
\begin{document}
\maketitle
\noindent\textit{Editorial note.} All delivered mathematical text is the author's own, copied from the pinned Stacks Project chapter file algebra.tex; no statement or proof step was rewritten, reordered or removed.\par\medskip
\noindent\textit{Editorial note.} Preamble provenance, as required for this source: the frozen snapshot stages the chapter's own preamble file beside it (retained-preamble.tex, 63 lines, 8116 bytes, the same shared Stacks-wide preamble the published chow and etale-cohomology cards use); the first line of the pinned chapter file is \textbackslash{}input\{preamble\}, so no preamble text lives inside the chapter itself. The card takes that file's lines 8-18 (the verbatim package and the three \textbackslash{}comment-based environments reference/slogan/history at lines 9-12, section 13.2, together with the Xy-pic setup) and its lines 154 to the runtime end of the file (the native theorem declarations and mathematical macros). Not one character was changed, no package was added, no macro was defined, and the excerpt serves only as preamble, never as a span coordinate.\par\medskip
\noindent\textit{Editorial note.} The same-chapter pointers are located at their \textbackslash{}label lines inside the pinned chapter file; the cross-chapter pointers are located in the project's own tag list staged in the snapshot (sources/stacks-chow/tags.txt), at the verbatim lines carrying their canonical tags.\par\medskip
\noindent\textit{Editorial note.} The counted spans contain no \textbackslash{}ref and no \textbackslash{}cite, so the card carries no navigation stub and no bibliography.\par\medskip
\noindent\textit{Allowed standard prerequisites (editorial).}\begin{itemize}
\item Flatness tested on quotients by powers of the ideal. Used in the first part of the counted proof. Reduces flatness of the module modulo each power to a Tor computation.
\item Long exact sequences of Tor for a short exact sequence. Used throughout the counted proof. Transports vanishing of Tor one from the quotient by the ideal to the powers and to modules annihilated by a power.
\end{itemize}\medskip
\begin{lemma}
\label{lemma-what-does-it-mean}
Let $R$ be a ring.
Let $I \subset R$ be an ideal.
Let $M$ be an $R$-module.
If $M/IM$ is flat over $R/I$ and $\text{Tor}_1^R(R/I, M) = 0$ then
\begin{enumerate}
\item $M/I^nM$ is flat over $R/I^n$ for all $n \geq 1$, and
\item for any module $N$ which is annihilated by $I^m$ for some $m \geq 0$
we have $\text{Tor}_1^R(N, M) = 0$.
\end{enumerate}
In particular, if $I$ is nilpotent, then $M$ is flat over $R$.
\end{lemma}

\begin{proof}
Assume $M/IM$ is flat over $R/I$ and $\text{Tor}_1^R(R/I, M) = 0$.
Let $N$ be an $R/I$-module. Choose a short exact sequence
$$
0 \to K \to \bigoplus\nolimits_{i \in I} R/I \to N \to 0
$$
By the long exact sequence of $\text{Tor}$ and the vanishing of
$\text{Tor}_1^R(R/I, M)$ we get
$$
0 \to \text{Tor}_1^R(N, M) \to K \otimes_R M \to
(\bigoplus\nolimits_{i \in I} R/I) \otimes_R M \to N \otimes_R M \to 0
$$
But since $K$, $\bigoplus_{i \in I} R/I$, and $N$ are all annihilated
by $I$ we see that
\begin{align*}
K \otimes_R M & = K \otimes_{R/I} M/IM, \\
(\bigoplus\nolimits_{i \in I} R/I) \otimes_R M & =
(\bigoplus\nolimits_{i \in I} R/I) \otimes_{R/I} M/IM, \\
N \otimes_R M & = N \otimes_{R/I} M/IM.
\end{align*}
As $M/IM$ is flat over $R/I$ we conclude that
$$
0 \to K \otimes_{R/I} M/IM \to
(\bigoplus\nolimits_{i \in I} R/I) \otimes_{R/I} M/IM \to
N \otimes_{R/} M/IM \to 0
$$
is exact. Combining this with the above we conclude that
$\text{Tor}_1^R(N, M) = 0$ for any $R$-module $N$ annihilated by $I$.

\medskip\noindent
Let us prove (2) by induction on $m$. The case $m = 1$ was
done in the previous paragraph. For $N$ annihilated by $I^m$
for $m > 1$ we may choose an exact sequence
$0 \to N' \to N \to N'' \to 0$ with $N'$ and $N''$ annihilated
by $I^{m - 1}$. For example one can take $N' = IN$ and $N'' = N/IN$.
Then the exact sequence
$$
\text{Tor}_1^R(N', M) \to
\text{Tor}_1^R(N, M) \to
\text{Tor}_1^R(N'', M)
$$
and induction prove the vanishing we want.

\medskip\noindent
Finally, we prove (1). Given $n \geq 1$ we have to show that $M/I^nM$
is flat over $R/I^n$. In other words, we have to show that the functor
$N \mapsto N \otimes_{R/I^n} M/I^nM$ is exact on the category of $R$-modules
$N$ annihilated by $I^n$. However, for such $N$ we have
$N \otimes_{R/I^n} M/I^nM = N \otimes_R M$. By the vanishing of
$\text{Tor}_1$ in (2) we see that the functor $N \mapsto N \otimes_R M$
is exact on the category of $N$ annihilated by some power of $I$
and we conclude.
\end{proof}

\end{document}
